\section*{Chapitre \ref{ch:g12:equilibrium} --- L'équilibre chimique~: quotient et constante}

\begin{solution}{exo:g12:equilibrium:1}
(a) $Q = \dfrac{[\ce{NH3}]^2 (c^\circ)^2}{[\ce{N2}][\ce{H2}]^3}$;
(b) $Q = \dfrac{[\ce{CH3COO-}][\ce{H3O+}]}{[\ce{CH3COOH}]\,c^\circ}$ (l'eau,
qui est le \omterm{def:g6:solutions-and-solubility:solute}{solvant}, est omise)~;
(c) $Q = [\ce{Ca^{2+}}][\ce{CO3^{2-}}]/(c^\circ)^2$ (le solide est
omis)~;
(d) $Q = \dfrac{[\ce{FeSCN^{2+}}]\,c^\circ}{[\ce{Fe^{3+}}][\ce{SCN-}]}$.
\end{solution}

\begin{solution}{exo:g12:equilibrium:2}
$\tau = 0.020/0.050 = 0.40$~: la réaction n'est pas totale.
\end{solution}

\begin{solution}{exo:g12:equilibrium:3}
$Q = 2 < K$~: sens direct. $Q = 50 > K$~: sens inverse. $Q = K$~: pas
d'évolution.
\end{solution}

\begin{solution}{exo:g12:equilibrium:4}
La composition est constante, mais les réactions dans les sens direct et
inverse se poursuivent à la même vitesse~: des \omterm{def:g7:atoms-and-molecules:molecule}{molécules} continuent de
réagir dans les deux sens.
\end{solution}

\begin{solution}{exo:g12:equilibrium:5}
Non, $K$ n'est pas modifiée (l'\omterm{def:g10:reaction-progress-table:system}{état final} est le même)~; la durée
nécessaire pour atteindre l'équilibre est plus courte.
\end{solution}

\begin{solution}{exo:g12:equilibrium:6}
$\dfrac{x^2}{(2-x)^2} = 4$, donc $\dfrac{x}{2-x} = 2$ et
$x = \qty{1.33}{mol}$~: \qty{0.67}{mol} d'acide et d'\omterm{def:g11:functional-groups:alcohol}{alcool},
\qty{1.33}{mol} d'\omterm{def:g11:functional-groups:carboxylic-acid}{ester} et d'eau (de nouveau les deux tiers).
\end{solution}

\begin{solution}{exo:g12:equilibrium:7}
$Q = (0.5 \times 0.5)/(0.5 \times 0.5) = 1 < 4$~: sens direct, il se
forme davantage d'\omterm{def:g11:functional-groups:carboxylic-acid}{ester}.
\end{solution}

\begin{solution}{exo:g12:equilibrium:8}
Environ \qty{0.52}{mol} du côté de l'acide et \qty{0.79}{mol} du côté de
l'\omterm{def:g11:functional-groups:carboxylic-acid}{ester}. Au bout d'environ \qty{100}{h}, les deux courbes rejoignent la
ligne en tirets~: c'est l'équilibre.
\end{solution}

\begin{solution}{exo:g12:equilibrium:9}
Avec $y$ mol hydrolysées~: $\dfrac{y^2}{(1-y)^2} = \dfrac{1}{4}$, donc
$\dfrac{y}{1-y} = \dfrac{1}{2}$, $y = \frac{1}{3}$~: $\frac{2}{3}$ mol
d'\omterm{def:g11:functional-groups:carboxylic-acid}{ester} et d'eau, $\frac{1}{3}$ mol d'acide et d'\omterm{def:g11:functional-groups:alcohol}{alcool}. C'est le même
\omterm{def:g6:pure-substances-and-mixtures:mixture}{mélange} final qu'avec l'estérification.
\end{solution}

\begin{solution}{exo:g12:equilibrium:10}
Un solide est une phase pure, dont la «~concentration~» ne change pas~:
il est omis dans $Q$. Ajouter du solide ne change pas $Q$, et ne déplace
donc pas l'équilibre (tant qu'il reste du solide).
\end{solution}

\begin{solution}{exo:g12:equilibrium:11}
$K$ fixe le rapport entre le dioxyde de carbone du gaz et celui de
l'eau. Une fois la bouteille ouverte, le gaz situé au-dessus de l'eau
s'échappe dans l'\omterm{def:g4:air-a-mixture-of-gases:air}{air}~: le dioxyde de carbone du gaz diminue, $Q < K$, et
davantage de gaz quitte l'eau, jusqu'à ce qu'il n'en reste presque plus.
\end{solution}

\begin{solution}{exo:g12:equilibrium:12}
À l'équilibre~: acide et \omterm{def:g11:functional-groups:alcohol}{alcool} $\frac{1}{3}$, \omterm{def:g11:functional-groups:carboxylic-acid}{ester} et eau
$\frac{2}{3}$. Retirer la moitié de l'eau en laisse $\frac{1}{3}$~:
$Q = \dfrac{(2/3)(1/3)}{(1/3)(1/3)} = 2 < 4$, sens direct. Si $y$ mol
de plus réagissent~: $(2/3 + y)(1/3 + y) = 4(1/3 - y)^2$, soit
$27y^2 - 33y + 2 = 0$, $y = 0.064$~: l'\omterm{def:g11:functional-groups:carboxylic-acid}{ester} passe à \qty{0.73}{mol}.
\end{solution}

\begin{solution}{exo:g12:equilibrium:13}
Le quotient de la réaction inverse a son numérateur et son dénominateur
échangés~: $Q' = 1/Q$, donc à l'équilibre $K' = 1/K$. Hydrolyse~:
$1/4 = 0.25$.
\end{solution}

\begin{solution}{exo:g12:equilibrium:14}
$Q = (2 \times 2)/(1 \times 1) = 4 = K$~: le \omterm{def:g6:pure-substances-and-mixtures:mixture}{mélange} est déjà à
l'équilibre~; sa composition ne change pas.
\end{solution}

\begin{solution}{exo:g12:equilibrium:15}
$n - 0.95 = 0.9025 / (4 \times 0.05) = 4.51$, donc $n = \qty{5.5}{mol}$
d'\omterm{def:g11:functional-groups:alcohol}{alcool} par \omterm{def:g10:the-mole:amount}{mole} d'acide.
\end{solution}

\begin{solution}{pb:g12:equilibrium:1}
\textbf{1.} \ce{CH3-COOH + CH3-CH2-OH <=> CH3-COO-CH2-CH3 + H2O}.

\textbf{2.} Elle n'est pas totale~: elle s'arrête dans un équilibre où
les deux sens se poursuivent.

\textbf{3.} $Q = \dfrac{n_{\text{ester}}\, n_{\text{eau}}}{n_{\text{acide}}\, n_{\text{alcool}}}$.

\textbf{4.} $x_{\max} = \qty{1}{mol}$~; $x_f = \frac{2}{3}$ mol.

\textbf{5.} $\tau = 0.67$.

\textbf{6.} Acide et \omterm{def:g11:functional-groups:alcohol}{alcool}~: $\frac{1}{3}$ mol chacun~; \omterm{def:g11:functional-groups:carboxylic-acid}{ester} et eau~:
$\frac{2}{3}$ mol chacun.

\textbf{7.} $K = \dfrac{(2/3)^2}{(1/3)^2} = 4$.

\textbf{8.} Un tiers hydrolysé~: \omterm{def:g11:functional-groups:carboxylic-acid}{ester} et eau $\frac{2}{3}$, acide et
\omterm{def:g11:functional-groups:alcohol}{alcool} $\frac{1}{3}$~: le même \omterm{def:g6:pure-substances-and-mixtures:mixture}{mélange}, $Q_{eq} = 4$.

\textbf{9.} C'est la même réaction à la même température, et $K$ ne
dépend que de la température.

\textbf{10.} Acide $1 - x$~; \omterm{def:g11:functional-groups:alcohol}{alcool} $3 - x$~; \omterm{def:g11:functional-groups:carboxylic-acid}{ester} $x$~; eau $x$.

\textbf{11.} $x^2 = 4(1 - x)(3 - x) = 4(3 - 4x + x^2)$, donc
$3x^2 - 16x + 12 = 0$.

\textbf{12.} $x = (16 \pm \sqrt{256 - 144})/6 = (16 \pm 10.58)/6$~:
0.903 ou 4.43~; seule $x_f = \qty{0.903}{mol}$ est comprise entre 0 et 1.

\textbf{13.} \qty{90}{\percent}.

\textbf{14.} $M = \qty{88.0}{g/mol}$: $0.903 \times 88.0 = \qty{79.5}{g}$.

\textbf{15.} Chaque élimination d'eau diminue le numérateur de $Q$, qui
passe sous $K$~: le système évolue de nouveau dans le sens direct.

\textbf{16.} La réaction se poursuivrait jusqu'à l'épuisement de
l'acide~: \qty{100}{\percent}.

\textbf{17.} L'éthanol est moins cher, et son excès se sépare facilement
de l'\omterm{def:g11:functional-groups:carboxylic-acid}{ester} (et se recycle).

\textbf{18.} Seulement la durée~: une réaction qui ne dégage presque pas
de chaleur a une constante qui varie à peine avec la température, et le
chauffage ne fait qu'accélérer l'arrivée à l'équilibre.

\textbf{19.} Environ \qty{90}{\percent} de l'acide.
\end{solution}
